% ======================================================================
% Supplementary Information
% Display items demoted from the main text for npj Digital Medicine.
% Figures renumber as S1, S2, ...; tables as ST1, ST2, ...
% All cross-references keep their original labels and auto-resolve.
% ======================================================================
\clearpage
\setcounter{figure}{0}
\setcounter{table}{0}
\renewcommand{\thefigure}{S\arabic{figure}}
\renewcommand{\thetable}{S\arabic{table}}

\section*{Supplementary Information}

% ----------------------------------------------------------------------
% Supplementary notes
% ----------------------------------------------------------------------

\subsection*{Supplementary Note 1: extended protocol for domains B and C}

\paragraph{Glioma seed and propagation.}
The VASARI-feature classifier is trained on the TCIA-annotated seed---$178$
LGG and $75$ GBM subjects carry expert scores, of which $63$ LGG also carry
the co-registered preoperative MRI and molecular call the protocol
needs---and propagated to the ${\sim}1{,}251$ BraTS 2021
subjects.\cite{baid2021brats} This occupies the position that \fuzzyarc{}\,V2,
trained on Zenodo-ANORAK and propagated to TCGA-LUAD, occupies in the LUAD
two-channel architecture (\cref{fig:two_channels}). Because the VASARI
channel on BraTS is therefore \emph{propagated} rather than expert-scored
(held-out feature $\auroc\!\lesssim\!0.6$), the MGMT claim in the main text
rests on the unconstrained encoder remaining at chance, not on the strict
empty-atom versus data-floor separation that real VASARI licenses for
\textit{KEAP1} and \textit{RBM10} in LUAD.

\paragraph{Glioma pilot, per condition.}
On the $63$-subject expert-scored cohort, five-fold patient-disjoint
(\cref{tab:glioma_domainb}): IDH reaches $\auroc\!=\!0.86$ on the encoder
alone with the fusion arms adding nothing beyond a fold's standard deviation,
the radiogenomic expectation for regime~(i). For 1p/19q neither channel alone
is informative ($\mathrm{B2}\!=\!0.60$, $\mathrm{B3}\!=\!0.50$) while their
concatenation recovers it ($\mathrm{PI}\!=\!0.74$), the signature of
complementary unique atoms rather than of a synergistic one, which places it
in sub-form (ii-a). The expectation that the non-additive aggregator would
win on 1p/19q was our modelling hypothesis, not an established radiogenomic
result; at $n\!\sim\!63$ the learned Choquet pair interactions shrink to
$\lvert I_{jk}\rvert\!\lesssim\!0.005$, so that read is deferred rather than
asserted. The transfer therefore did not weaken the regime-(ii) claim---which
stands on $\mathrm{PI}\!>\!\max(\mathrm{B2},\mathrm{B3})$---but declined the
stronger sub-claim at a sample size that cannot support it.

\paragraph{Electrocardiogram, measured and remaining targets.}
The $71$ SCP-ECG statements are organised into diagnostic, form and rhythm
categories with a proposed hierarchy,\cite{wagner2020ptbxl} and PTB-XL$+$
adds upward propagation through the SNOMED-CT
tree.\cite{strodthoff2023ptbxlplus} The measured pair reported in the main
text is sex ($\auroc\!=\!0.90$) and age$\,\geq\!65$ ($0.84$) under the best
aggregator, against $0.62$ and $0.73$ for the statement channel alone, with
$\mathrm{FC}-\mathrm{PI}\!=\!+0.002$ and $-0.002$ respectively
(\cref{tab:ecg_domainc}). The $59$ borderline laboratory indicators that
\citet{xiao2025anyecglab} describe as information that ``may exist in the
signal but is difficult to extract'' are the natural next targets. Beyond
them, structural heart disease labels from EchoNext\cite{hughes2026echonext}
and the paired ECG/chest-radiograph/laboratory records of
MIMIC-IV-ECG\cite{gow2023mimicivecg} extend the same protocol to a
multimodal setting in which the missingness correction of
\citet{choi2025icym2i} becomes mandatory rather than optional.

% ----------------------------------------------------------------------
% Supplementary figures
% ----------------------------------------------------------------------

% --- S1: six-level schematic stack ---
\begin{figure}[ht!]
  \centering
  \includegraphics[width=\linewidth]{lchai_six_level_stack}
  \caption{\textbf{\lchai{} v2 system overview and six-level
  interpretability stack.} Each WSI passes through tiling, the
  frozen \ctranspath{} foundation encoder, the ANORAK-trained
  pattern classifier (6-d fuzzy membership $\mathbf{p}\!\in\!\Delta^5$),
  and the \fcmil{} 2-additive Choquet aggregator. From a single
  forward pass through the aggregator the system materialises six
  interpretability layers: \textbf{L1} ABMIL attention,
  \textbf{L2} pattern overlay, \textbf{L3} SHAP channel
  decomposition, \textbf{L4} Choquet--Shapley values and pairwise
  interaction indices,   \textbf{L5} channel ablation, and \textbf{L6} fuzzy linguistic
  labels.
  The explanation layer wraps these six artefacts in an LLM agent
  grounded in a three-layer knowledge graph
  (NCIt+MONDO+curated PMID edges), with the fuzzy labels
  constraining \emph{how} the agent speaks and the KG constraining
  \emph{what} entities and edges may appear.}
  \label{fig:lchai_six_levels}
\end{figure}

% --- S2: per-gene case grid ---
\begin{figure}[ht!]
  \centering
  \begin{subfigure}[t]{0.32\linewidth}
    \casepanel{lchai_v2_TP53_TCGA-99-8025_combined}
    \caption{TP53, TCGA-99-8025 ($P=0.951$, TP).}
    \label{fig:case_tp53}
  \end{subfigure}\hfill
  \begin{subfigure}[t]{0.32\linewidth}
    \casepanel{lchai_v2_EGFR_TCGA-86-8280_combined}
    \caption{EGFR, TCGA-86-8280 (TP).}
    \label{fig:case_egfr}
  \end{subfigure}\hfill
  \begin{subfigure}[t]{0.32\linewidth}
    \casepanel{lchai_v2_KRAS_TCGA-55-7815_combined}
    \caption{KRAS, TCGA-55-7815 ($P=0.965$, IMA proxy).}
    \label{fig:case_kras}
  \end{subfigure}\\[1.2ex]
  \begin{subfigure}[t]{0.32\linewidth}
    \casepanel{lchai_v2_STK11_TCGA-49-AAR0_combined}
    \caption{STK11, TCGA-49-AAR0 (TP).}
    \label{fig:case_stk11}
  \end{subfigure}\hfill
  \begin{subfigure}[t]{0.32\linewidth}
    \casepanel{lchai_v2_KEAP1_TCGA-49-AAR9_combined}
    \caption{KEAP1, TCGA-49-AAR9 (TP).}
    \label{fig:case_keap1}
  \end{subfigure}\hfill
  \begin{subfigure}[t]{0.32\linewidth}
    \casepanel{lchai_v2_RBM10_TCGA-78-7148_combined}
    \caption{RBM10, TCGA-78-7148 (TP).}
    \label{fig:case_rbm10}
  \end{subfigure}
  \caption{\textbf{Per-gene case studies.} For each of the six
  target genes, one representative true-positive slide is shown as
  the combined \lchai{} overlay: the \ctranspath{} tissue thumbnail
  of the whole-slide image with the L1 ABMIL attention map and the
  L2 ANORAK pattern saliency superimposed, so that the warm regions
  localise where high attention and high pattern-saliency coincide---
  the evidence the aggregator weighed---over the green tissue field.
  The KRAS panel illustrates the value of decoupling the
  cohort-level verdict (Inconclusive at $0.609$ mean AUROC,
  \cref{tab:mutation_luad_resultssec}) from
  the per-slide explanation: TCGA-55-7815 is per-slide confident
  ($P=0.965$) and the L4 Choquet readout (\cref{fig:kras_proxy_b})
  localises the confidence in the lepidic$\times$solid
  co-activation. Per-slide additional examples and full SHAP
  channel decompositions appear in Extended Data.}
  \label{fig:case_grid}
\end{figure}

% --- S3: three-layer knowledge graph ---
\begin{figure}[ht!]
  \centering
  \includegraphics[width=\linewidth]{kg_three_layers_method}
  \caption{\textbf{Three-layer knowledge graph
  $\mathcal{K} = (V_C \cup V_D \cup V_L,\, E)$ used to ground the
  LLM explanation layer.}
  $V_C$: NCIt~\cite{sioutos2007ncit} concept nodes (genes,
  pathways, histological patterns).
  $V_D$: MONDO~\cite{mondo_shefchek2020} disease nodes
  (e.g.\ \emph{lung adenocarcinoma},
  \emph{invasive mucinous adenocarcinoma}).
  $V_L$: a per-case dynamic literature layer built by DeepSearch
  at query time, each node carrying a PMID provenance tag that is
  exposed to the user as a click-through citation.
  The agent's prompt grammar permits it to invoke only entities
  and edges in $V$ and $E$; this constrains \emph{what} the
  narrative can claim. \cref{eq:trap_mf} constrains \emph{how}
  evidence strength is verbalised.}
  \label{fig:kg_three_layers}
\end{figure}

% --- S4: real ECG example ---
\begin{figure}[t]
\centering
\includegraphics[width=\linewidth]{ecg_ptbxl_example}
\caption{\textbf{A real 12-lead electrocardiogram from PTB-XL}
(a $37$-year-old woman, sinus rhythm, reported normal ECG), shown in the
standard clinical $3\times4$ format with a lead-II rhythm strip on an ECG
grid ($25$\,mm/s, $10$\,mm/mV). This is the raw $100$\,Hz signal the
encoder $\mathbf{E}$ ingests---the electrocardiographic counterpart of the
whole-slide histopathology images in LUAD. Patient age and sex are
recorded for every PTB-XL case but are absent from the $71$ SCP-ECG
diagnostic statements that form the source channel $\mathbf{P}$, so the
framework recovers both from exactly this signal
(\cref{tab:ecg_domainc,fig:pipeline_ecg}), a leakage-free
physiological-state translation.\cite{wagner2020ptbxl}}
\label{fig:ecg_ptbxl}
\end{figure}

% --- S5: real glioma MRI example ---
\begin{figure}[t]
\centering
\includegraphics[width=\linewidth]{glioma_example}
\caption{\textbf{A real preoperative multiparametric MRI from a TCGA-LGG
glioma} (an IDH-mutant, 1p/19q-codeleted case)---the four co-registered
sequences (T1, T1-Gd, T2, FLAIR) the encoder $\mathbf{E}$ ingests, shown at
the tumour-centroid axial slice with the expert tumour segmentation contour
(red). Below, the neuroradiologist VASARI scores
for the same case---the source channel $\mathbf{P}$ the framework fuses---are
the radiological counterpart of the ANORAK growth patterns in LUAD and the
SCP-ECG statements in the electrocardiogram
(\cref{tab:glioma_domainb,fig:pipeline_glioma}). This is the raw signal from
which IDH and 1p/19q are read, the imaging analogue of the whole-slide
histopathology in LUAD and the twelve-lead trace in
\cref{fig:ecg_ptbxl}.}
\label{fig:glioma_example}
\end{figure}

% ----------------------------------------------------------------------
% Supplementary tables
% ----------------------------------------------------------------------

% --- ST1: CPTAC power ---
\begin{table}[ht!]
\centering
% Table tab:cptac_power — expected positive-set sizes for the six target
% genes in the CPTAC-LUAD cohort (n=110), used to make the within-domain
% cohort-transfer pre-registration honest about which targets it can and
% cannot estimate. Prevalences are those reported for the CPTAC-LUAD
% discovery cohort; n_+ is the implied positive count at n=110. The
% estimation floor is the n_+=37 at which RBM10 is only underpowered in the
% larger TCGA-LUAD cohort of this paper.
\footnotesize
\setlength{\tabcolsep}{6pt}
\renewcommand{\arraystretch}{1.15}
\begin{tabular}{l c c c}
\toprule
\textbf{Gene} & \textbf{Regime (LUAD)} &
  \textbf{CPTAC prevalence} & \textbf{$n_+$ at $n\!=\!110$ (floor $37$)} \\
\midrule
TP53  & (i) saturation & $46\%$ & $\sim$51 \quad above \\
KRAS  & (ii) synergy   & $33\%$ & $\sim$36 \quad borderline \\
STK11 & (iii)          & $17\%$ & $\sim$19 \quad below \\
KEAP1 & (iii)          & $17\%$ & $\sim$19 \quad below \\
EGFR  & (i) saturation & $14\%$ & $\sim$15 \quad below \\
RBM10 & estimation floor & $\sim$5--8\%$^{\dagger}$ & $\sim$6--9 \quad below \\
\bottomrule
\end{tabular}
\begin{flushleft}\footnotesize
Prevalences are those reported for the CPTAC-LUAD discovery cohort
($n\!=\!110$)\cite{gillette2020cptacluad}; $^{\dagger}$\textit{RBM10} was
not tabulated with a separate frequency in that report and is estimated
from the TCGA-LUAD range (here $5.4\%$). $n_+$ is the implied positive
count; the floor is the $n_+\!=\!37$ at which \emph{RBM10} is only
underpowered in the $n\!=\!687$ TCGA-LUAD cohort reported here. Only
\emph{TP53} clears the floor and \emph{KRAS} reaches it; the four
remaining targets---including all of the regime-(iii) genes---fall
below it, so a cohort-transfer replication on CPTAC-LUAD can confirm or
refute the well-powered assignments (\emph{TP53}, \emph{KRAS}) with
overlapping confidence intervals but is, by construction, underpowered
for exactly the untranslatable candidates it would most need to check.
\end{flushleft}

\caption{Observed positive-set sizes for the six target genes in the
CPTAC-LUAD cohort-transfer test, from the actual non-silent mutation
frequencies in cBioPortal \texttt{luad\_cptac\_2020} ($n\!=\!110$). These
fix, in advance of the imaging read-out, which targets the transfer can
estimate: in the genotyped-with-imaging subset ($n\!=\!106$) only
\textit{TP53} clears the $n_+\!=\!37$ floor, \textit{EGFR} falls just
short, and the regime-(iii) genes and \textit{RBM10} fall well below.}
\label{tab:cptac_power}
\end{table}

% --- ST2: CPTAC-internal refit of the four conditions ---
\begin{table}[ht!]
\centering
% Table tab:cptac_refit — CPTAC-internal refit of the four conditions.
% Only the representation is frozen (CTransPath embeddings + ANORAK pattern
% probabilities recomputed on CPTAC tiles); the four aggregation conditions are
% re-fitted with 5-fold patient-disjoint CV inside CPTAC. TP53 only: it is the
% one gene clearing the n_+=37 estimation floor in this subset (n_+=55).
% "Native" tiles CPTAC at level 0 exactly as the TCGA pipeline does. "FOV-matched"
% centre-crops and rescales every tile so its physical field of view matches a
% 40x TCGA tile (target 0.25 micron/pixel), testing the scanner-magnification
% account of the transfer failure.
\footnotesize
\setlength{\tabcolsep}{6pt}
\renewcommand{\arraystretch}{1.2}
\begin{tabular}{l c c}
\toprule
\textbf{Condition (\textit{TP53}, $n_+\!=\!55$)} & \textbf{Native} & \textbf{FOV-matched} \\
\midrule
B2 \quad encoder only                & $0.665$ & $0.561$ \\
B3 \quad pattern channel only        & $0.667$ & $0.570$ \\
\piabmil{} \quad concat fusion       & $0.698$ & $0.632$ \\
\fcmil{} \quad $2$-additive Choquet  & $0.669$ & $0.538$ \\
\midrule
$\piabmil{}-\mathrm{B2}$             & $+0.032$ & $+0.073$ \\
Regime, point estimate               & (ii-a) & (ii-a) \\
Regime, adjudicated                  & unresolved & unresolved \\
$P(\text{regime reproduces TCGA})$   & $0.214$ & $0.129$ \\
\midrule
\multicolumn{3}{l}{\emph{Reference:} TCGA-LUAD $0.717$; frozen transfer $0.53$
($0.42$--$0.64$), \cref{tab:cptac_transfer}} \\
\bottomrule
\end{tabular}
\begin{flushleft}\footnotesize
Pooled out-of-fold AUROC over $5$ patient-disjoint folds on the $106$
genotyped CPTAC patients ($498$ slides). Only the \emph{representation} is
carried over from TCGA; the aggregation conditions are re-fitted inside CPTAC,
which is the comparison the frozen-pipeline test of \cref{tab:cptac_transfer}
cannot make. The native arm recovers $0.698$ against the TCGA value of $0.717$,
so the morphology-to-genotype coupling for \textit{TP53} is present in CPTAC at
comparable strength and the off-the-shelf collapse to $0.53$ is a property of
the transported decision function, not of the encoder or of the tissue.
The \textbf{FOV-matched} arm is a negative control for the
scanner-magnification account (TCGA-LUAD is $76.4\%$ $40\times$, CPTAC-LUAD
$98.6\%$ $20\times$): matching the physical field of view lowers every
condition and inflates fold-to-fold spread (B2 ranges $0.373$--$0.918$ across
folds, against $0.587$--$0.791$ natively), so scale mismatch is not the
mechanism, though the centre-crop is itself lossy and the arm is not evidence
in the opposite direction either.
Regime adjudication follows the noise-aware rule of the Methods: the deciding
contrast is $\piabmil{}-\mathrm{B2}$, whose $2000$-sample paired-bootstrap
$95\%$ CI is $-0.041$ to $+0.102$ natively and $-0.038$ to $+0.174$
FOV-matched. Both span zero, so the point-estimate read of (ii-a) is
\emph{not} separable from noise and no regime is asserted; the cohort neither
confirms nor contradicts the TCGA reading of~(i).
\end{flushleft}

\caption{\textbf{Re-fitting the four conditions inside CPTAC-LUAD, with
only the representation carried over from TCGA.} This is the comparison
the frozen-pipeline transfer of \cref{tab:cptac_transfer} cannot make,
and it separates the two readings that transfer alone leaves open. The
recovered values, the field-of-view control and the regime adjudication
are set out below the table.}
\label{tab:cptac_refit}
\end{table}

% --- ST3: mutation prevalence ---
\begin{table}[ht!]
\centering
% Table tab:mutation_prevalence — TCGA-LUAD cohort mutation prevalence
% (687 analysable diagnostic slides, 668 patients).
\small
\setlength{\tabcolsep}{8pt}
\begin{tabular}{l c c c}
\toprule
\textbf{Gene} & \textbf{Positive slides} ($n_{+}$) &
  \textbf{Prevalence} (\%) & \textbf{Conclusive verdict?} \\
\midrule
\textit{TP53}  & 261 & 38.0 & Yes \\
\textit{KRAS}  & 196 & 28.5 & No  \\
\textit{KEAP1} &  98 & 14.3 & No  \\
\textit{EGFR}  &  97 & 14.1 & Yes \\
\textit{STK11} &  68 &  9.9 & No  \\
\textit{RBM10} &  37 &  5.4 & No  \\
\bottomrule
\end{tabular}
\begin{flushleft}\footnotesize
$n_{+}\!=\!$~number of slides with at least one non-silent somatic
mutation called for the gene under the TCGA harmonised pipeline.
\emph{Conclusive verdict} is the cohort-level signal-presence flag
(mean AUROC $\geq\!0.70$ under five-fold cross-validation; see
Methods, Cross-validation and statistical analysis).
\end{flushleft}

\caption{Mutation prevalence in the analysable TCGA-LUAD cohort
($n\!=\!687$ slides from $668$ patients).
The Conclusive-verdict column anticipates the cohort-level outcome
of the five-fold cross-validation reported in Results.}
\label{tab:mutation_prevalence}
\end{table}

% --- ST4: verdict bands ---
\begin{table}[ht!]
\centering
% Table tab:viz_folds_methods — Conclusive / Inconclusive verdict bands
% over the five-fold patient-stratified CV protocol.
% Wide table (4 cols, last col long): use tabularx so the "Genes" column
% wraps cleanly inside the text width.
\small
\setlength{\tabcolsep}{6pt}
\begin{tabularx}{\linewidth}{l c c X}
\toprule
\textbf{Verdict} & \textbf{Mean AUROC band} &
  \makecell{\textbf{Per-slide score}\\\textbf{still reported?}} &
  \textbf{Genes (this cohort)} \\
\midrule
Conclusive   & $\geq 0.70$ & Yes &
  \textit{TP53} ($0.717$), \textit{EGFR} ($0.701$)$^{\dagger}$ \\
Inconclusive & $< 0.70$    & Yes &
  \textit{KRAS} ($0.609$), \textit{STK11} ($0.695$)$^{\dagger}$,
  \textit{KEAP1} ($0.610$), \textit{RBM10} ($0.661$) \\
\bottomrule
\end{tabularx}
\begin{flushleft}\footnotesize
The verdict is a population-level signal-presence flag, not a
per-slide confidence threshold; the per-slide model probability is
reported independently in every case.
$^{\dagger}$EGFR and STK11 are within $\pm 0.01$ of the threshold
under the upper-tier conditions and are reported here under their
best-tier AUROC. The Conclusive/Inconclusive flag is robust to
choice of MIL aggregator within the upper tier; the per-gene
ranking is not.
\end{flushleft}

\caption{Conclusive/Inconclusive verdict bands and their assignment
to the six target genes under the five-fold patient-stratified CV
protocol used throughout the paper.}
\label{tab:viz_folds_methods}
\end{table}
